\documentclass[a4paper,14pt]{extarticle}

\usepackage[T2A]{fontenc}
\usepackage[utf8]{inputenc}
\usepackage[russian]{babel}
\usepackage{geometry}
\geometry{left=3cm,right=1.5cm,top=2cm,bottom=2cm}
\usepackage{setspace}
\onehalfspacing
\usepackage{indentfirst}
\usepackage{tempora}
\pagestyle{empty}
\setlength{\parindent}{1.25cm}
\setlength{\parskip}{0pt}

\begin{document}

\begin{center}
\begin{minipage}{\linewidth}
\centering\bfseries
Отзыв научного руководителя\\
на диссертацию Горбова Григория Владиславовича\\
«Методы и алгоритмы безопасного обучения с подкреплением в робототехнических задачах», представленную на соискание учёной степени кандидата технических наук по специальности 1.2.1~--- «Искусственный интеллект и машинное обучение».
\end{minipage}
\end{center}

Разработка методов безопасного обучения с подкреплением является актуальным направлением в области искусственного интеллекта и робототехники. Современные автономные агенты должны не только достигать поставленных целей, но и соблюдать ограничения безопасности при взаимодействии со сложной и частично неизвестной средой. Эта задача особенно важна для робототехнических систем, поскольку ошибки в процессе обучения и исследования среды могут приводить к столкновениям и наносить вред человеку. Всё это определяет актуальность диссертационной работы Григория.

В первой главе проведён обзор существующих методов безопасного обучения с подкреплением, а также сред для их оценки. Методы были разделены на группы, и выявлены основные проблемы каждой из них: подходы оптимизации стратегии с ограничениями не гарантируют безопасность каждого отдельного действия; методы на основе фильтрации действий требуют априорных знаний о среде; алгоритмы безопасного обучения на основе моделей мира плохо оценивают потенциальное нарушение ограничений; методы безопасного иерархического обучения с подкреплением не гарантируют безопасность выбираемых подцелей, а также рассматриваются в рамках задач с коротким горизонтом.

Во второй главе предложена функция оценки будущих накопленных штрафов в рамках подхода SPOWL для улучшения скрытых представлений модели мира, а также для более точной оценки безопасности траекторий при планировании. Дополнительно предложен способ построения и агрегации функции оценки будущих накопленных штрафов. Эксперименты на средах Safety Gym показали лучшую производительность и безопасность по сравнению с предыдущими методами безопасного обучения с подкреплением. Отдельно показан вклад использования глобальной оценки безопасности траектории по сравнению с локальной оценкой.

В третьей главе была показана проблема ухудшения способностей к исследованию у методов безопасного обучения с подкреплением, и был предложен подход SG-Safe, использующий иерархическую структуру для решения этой проблемы. Метод использует генерацию промежуточных подцелей как дополнительный обучающий сигнал для безопасной стратегии. Такой подход позволяет агенту активнее исследовать среду и находить траектории к цели, не сводя поведение к чрезмерно осторожному избеганию любых потенциально опасных действий. Метод показывает лучшее качество в среде POLAMP в сравнении с методами безопасного обучения с подкреплением и методами на основе классического планирования.

В четвертой главе был предложен метод ITES, в котором безопасность учитывается как на уровне выбора подцелей, так и на уровне их достижения. Для обеспечения безопасности подцелей используются модель и критик штрафов, для обеспечения безопасного достижения подцелей стратегией применяется модель мира. Метод достигает лучшего качества в задачах навигации и манипуляции с большим горизонтом в сравнении с методами на основе моделей мира в безопасном обучении с подкреплением.

В пятой главе представлен бенчмарк Caged CrafText, предназначенный для обучения и оценки безопасного обучения с подкреплением в задачах, где цель и ограничения безопасности задаются на естественном языке. Бенчмарк включает различные типы ограничений и позволяет анализировать способность агента понимать инструкцию, связывать её с объектами и состояниями среды, планировать поведение и соблюдать требования безопасности, а также оценивать способности агента: понимание языка, связывание инструкций и текстовых ограничений с объектами из среды, память, безопасное исследование. Эксперименты показывают, что существующие методы безопасного обучения с подкреплением плохо справляются с ограничениями и задачами, заданными на естественном языке, в основном из-за консервативного поведения и отсутствия навыка связывания задачи и наблюдения.

В ходе диссертационного исследования Григорий показал высокий уровень самостоятельности и профессиональной подготовки. Он продемонстрировал способность формулировать научные задачи, разрабатывать алгоритмические решения, реализовывать экспериментальные системы и анализировать полученные результаты. Он также продемонстрировал способности к работе в команде при выполнении сложных научно-технических задач.

Основные результаты диссертационной работы опубликованы в рецензируемых научных изданиях и материалах международной конференции — в 5 публикациях: 2~--- в журналах категории К1 собственного перечня журналов МФТИ, 2~--- в сборниках трудов конференций уровня A и A*, 1~--- в виде препринта на arXiv.

Считаю, что данная диссертационная работа удовлетворяет всем требованиям Положения МФТИ о порядке присуждения учёных степеней, а её автор, Горбов Григорий Владиславович, заслуживает присуждения учёной степени кандидата технических наук по специальности 1.2.1~--- «Искусственный интеллект и машинное обучение».

\vspace{1\baselineskip}
\noindent Научный руководитель,\\
доктор физико-математических наук, профессор РАН\\
Панов Александр Игоревич

\vspace{2\baselineskip}
\noindent «\underline{\hspace{1cm}}» \underline{\hspace{3cm}} 2026~г.

\end{document}
